\chapter{Marco macroeconómico}

\marginal{el ajuste llega por el gasto}
El Paquete Económico 2025 es el primero de la administración entrante y el primero
en cinco años que propone una consolidación. Los requerimientos financieros del
sector público pasan de 5.9\% del PIB estimado para 2024 a \textbf{3.9\% en 2025},
y el balance primario cambia de signo: de un déficit de 1.4 puntos a un superávit
de 0.6. Dos puntos del PIB de ajuste en un ejercicio.

El marco macroeconómico que lo sostiene es menos optimista que los de los años
previos, y en un punto concreto es más conservador que el propio límite legal que
la Secretaría se impone.

\begin{table}[htbp]
\centering
\caption{Marco macroeconómico, 2024--2025}
\label{tab:C1_1}
\small
\setlength{\tabcolsep}{4pt}
\begin{tabular}{L{6.0cm}ccc}
\toprule
{Variable} & {2024 aprobado} & {2024 estimado} & {2025} \\
\midrule
{PIB, crecimiento real (rango)} & {[2.5,3.5]} & {[1.5,2.5]} & {[2.0,3.0]} \\
{PIB nominal (miles de millones de pesos)} & \num{34374} & \num{33927.7} & \num{36166.4} \\
{Deflactor del PIB (\%)} & \num{4.8} & \num{4.6} & \num{4.3} \\
{Inflación dic/dic (\%)} & \num{3.8} & \num{4.3} & \num{3.5} \\
{Inflación promedio (\%)} & \num{4.5} & \num{4.7} & \num{3.8} \\
{Tipo de cambio fin de periodo} & \num{17.6} & \num{19.7} & \num{18.5} \\
{Tipo de cambio promedio} & \num{17.1} & \num{18.2} & \num{18.7} \\
{Cetes 28 nominal fin de periodo (\%)} & \num{9.5} & \num{10.0} & \num{8.0} \\
{Cetes 28 nominal promedio (\%)} & \num{10.3} & \num{10.7} & \num{8.9} \\
{Cetes 28 real acumulada (\%)} & \num{6.7} & \num{6.7} & \num{5.6} \\
{Cuenta corriente (millones de dólares)} & \num{-14954} & \num{-7097.5} & \num{-7941.0} \\
{Cuenta corriente (\% del PIB)} & \num{-0.7} & \num{-0.4} & \num{-0.4} \\
{PIB EE.UU., crecimiento real (\%)} & \num{1.8} & \num{2.7} & \num{2.2} \\
{Producción industrial EE.UU. (\%)} & \num{2.0} & \num{0.5} & \num{2.0} \\
{Inflación EE.UU. promedio (\%)} & \num{2.4} & \num{2.9} & \num{2.2} \\
{SOFR 3 meses promedio (\%)} & \num{4.3} & \num{5.0} & \num{3.3} \\
{Fed Funds Rate promedio (\%)} & \num{4.5} & \num{5.3} & \num{3.8} \\
{Petróleo, precio promedio (dls/barril)} & \num{56.7} & \num{70.7} & \num{57.8} \\
{Plataforma de producción (mbd)} & \num{1983} & \num{1877} & \num{1891} \\
{Plataforma de exportación (mbd)} & \num{994} & \num{900} & \num{892} \\
{Plataforma de privados (mbd)} & \num{86} & \num{67} & \num{74} \\
{Gas, precio promedio (dls/MMBtu)} & \num{3.5} & \num{2.2} & \num{3.1} \\
\bottomrule
\end{tabular}
\par\vspace{2pt}\footnotesize Fuente: elaboración propia del ITED con información de la SHCP, CGPE 2025, Anexo II.5, p. 81. El CGPE 2025 se publicó sin capa de texto; las cifras se leyeron de la página renderizada y se validaron por identidad contable.
\end{table}


\section{Crecimiento y potencial}

La Secretaría prevé para 2025 un crecimiento real en el rango de 2.0 a 3.0\%,
frente a un rango estimado de 1.5 a 2.5 para el cierre de 2024. Pero la cifra que
importa para las finanzas públicas no es el rango sino la que está implícita en el
PIB nominal: con un PIB de 33,927.7 miles de millones en 2024 y 36,166.4 en 2025,
y un deflactor de 4.3\%, \textbf{el crecimiento real implícito es de 2.21\%}, por
debajo del punto medio del rango declarado.

\marginal{por debajo del potencial}
El crecimiento potencial que la propia Secretaría calcula para 2025, con la
metodología del reglamento de la ley de presupuesto, es de \textbf{2.32\%}: el
promedio de una tasa histórica de 2.19 y una prospectiva de 2.44. El crecimiento
implícito en el marco queda por debajo de él. Es un dato con consecuencia legal
además de analítica, porque la regla de deuda se activa según esa comparación, y
conviene registrar que el paquete 2025 se planta del lado conservador de una regla
que en ejercicios anteriores se leyó del lado contrario.

\section{Precios, tasas y tipo de cambio}

La inflación proyectada baja de 4.3\% en el cierre estimado de 2024 a 3.5\% en
diciembre de 2025, y el promedio de 4.7 a 3.8. La tasa nominal de referencia cae
con ella, de 10.7 a 8.9\% en promedio, de modo que \textbf{la tasa real promedio
baja de 5.7 a 4.9\%}. La distinción importa y este documento vuelve sobre ella en
el capítulo de deuda: bajar la inflación con la tasa nominal cayendo más deprisa
reduce la tasa real, y es la tasa real la que mueve el acervo.

El tipo de cambio es el supuesto que más llama la atención. La Secretaría proyecta
una \textbf{apreciación} del peso, de 19.7 a 18.5 por dólar al cierre del año, al
tiempo que el promedio anual se deprecia de 18.2 a 18.7. Los dos movimientos son
compatibles porque describen momentos distintos, pero conviene decir cuál se usa
en cada cálculo: el efecto sobre el acervo de deuda se calcula con el de fin de
periodo, y ahí la apreciación proyectada vale casi ocho décimas de punto del PIB.

\section{Petróleo}

El precio de la mezcla mexicana pasa de 70.7 dólares por barril estimados en 2024
a \textbf{57.8 en 2025}, una caída de 18\%. La plataforma de producción sube
levemente, de 1,877 a 1,891 miles de barriles diarios, mientras la de exportación
baja de 900 a 892. La combinación de menor precio con exportación estable es la
que explica por qué los ingresos petroleros del capítulo 3 crecen poco pese a que
el sector no se contrae.

\section{Mediano plazo}

\begin{table}[htbp]
\centering
\caption{Marco macroeconómico de mediano plazo, 2023--2030}
\label{tab:C1_2}
\footnotesize
\setlength{\tabcolsep}{4pt}
\resizebox{\textwidth}{!}{%
\begin{tabular}{L{4.6cm}cccccccc}
\toprule
{Variable} & {2023} & {2024} & {2025} & {2026} & {2027} & {2028} & {2029} & {2030} \\
\midrule
{PIB, crecimiento real (rango)} & \num{3.2} & {[1.5,2.5]} & {[2.0,3.0]} & {[2.0,3.0]} & {[2.0,3.0]} & {[2.0,3.0]} & {[2.0,3.0]} & {[2.0,3.0]} \\
{PIB nominal (miles de millones de pesos)} & \num{31772.4} & \num{33927.7} & \num{36166.4} & \num{38348.3} & \num{40682.5} & \num{43152.6} & \num{45789.6} & \num{48574.0} \\
{Deflactor del PIB promedio (\%)} & \num{4.5} & \num{4.6} & \num{4.3} & \num{3.5} & \num{3.5} & \num{3.5} & \num{3.5} & \num{3.5} \\
{Inflación dic/dic (\%)} & \num{4.7} & \num{4.3} & \num{3.5} & \num{3.0} & \num{3.0} & \num{3.0} & \num{3.0} & \num{3.0} \\
{Inflación promedio (\%)} & \num{5.5} & \num{4.7} & \num{3.8} & \num{3.2} & \num{3.0} & \num{3.0} & \num{3.0} & \num{3.0} \\
{Tipo de cambio fin de periodo} & \num{16.9} & \num{19.7} & \num{18.5} & \num{18.0} & \num{18.1} & \num{18.3} & \num{18.5} & \num{18.7} \\
{Tipo de cambio promedio} & \num{17.8} & \num{18.2} & \num{18.7} & \num{18.5} & \num{18.7} & \num{18.9} & \num{19.1} & \num{19.3} \\
{Cetes 28 nominal fin de periodo (\%)} & \num{11.3} & \num{10.0} & \num{8.0} & \num{7.0} & \num{5.5} & \num{5.5} & \num{5.5} & \num{5.5} \\
{Cetes 28 nominal promedio (\%)} & \num{11.1} & \num{10.7} & \num{8.9} & \num{7.4} & \num{6.2} & \num{5.5} & \num{5.5} & \num{5.5} \\
{Cetes 28 real acumulada (\%)} & \num{6.7} & \num{6.7} & \num{5.6} & \num{4.6} & \num{3.2} & \num{2.6} & \num{2.6} & \num{2.6} \\
{Cetes 28 real fin de periodo (\%)} & \num{6.3} & \num{5.4} & \num{4.4} & \num{3.9} & \num{2.4} & \num{2.4} & \num{2.4} & \num{2.4} \\
{Cetes 28 real promedio (\%)} & \num{5.3} & \num{5.7} & \num{4.9} & \num{4.1} & \num{3.1} & \num{2.4} & \num{2.4} & \num{2.4} \\
{Cuenta corriente (millones de dólares)} & \num{-5476.7} & \num{-7097.5} & \num{-7941.0} & \num{-10359.3} & \num{-10853.3} & \num{-11395.5} & \num{-11974.6} & \num{-12579.2} \\
{Cuenta corriente (\% del PIB)} & \num{-0.3} & \num{-0.4} & \num{-0.4} & \num{-0.5} & \num{-0.5} & \num{-0.5} & \num{-0.5} & \num{-0.5} \\
{Petróleo, precio promedio (dls/barril)} & \num{70.9} & \num{70.7} & \num{57.8} & \num{61.7} & \num{60.9} & \num{60.5} & \num{60.2} & \num{60.0} \\
{Plataforma de producción (mbd)} & \num{1942} & \num{1877} & \num{1891} & \num{1902} & \num{1911} & \num{1921} & \num{1930} & \num{1941} \\
{Plataforma de exportación (mbd)} & \num{1032} & \num{900} & \num{892} & \num{883} & \num{875} & \num{866} & \num{858} & \num{850} \\
{Gas, precio promedio (dls/MMBtu)} & \num{2.5} & \num{2.2} & \num{3.1} & \num{3.6} & \num{3.6} & \num{3.6} & \num{3.4} & \num{3.4} \\
\bottomrule
\end{tabular}
}
\par\vspace{2pt}\footnotesize Fuente: elaboración propia del ITED con información de la SHCP, CGPE 2025, Anexo III.1, p. 84. El CGPE 2025 se publicó sin capa de texto; las cifras se leyeron de la página renderizada y se validaron por identidad contable.
\end{table}


El escenario de mediano plazo supone crecimiento de 2.0 a 3.0\% todos los años
hasta 2030, inflación de 3.0 desde 2027, y una tasa nominal que baja hasta 5.5\%.
Es un escenario de convergencia sin sobresaltos. \textbf{Su rasgo más fuerte es lo
que hace con la deuda}: el saldo histórico de los requerimientos financieros se
mantiene exactamente en 51.4\% del PIB de 2025 a 2030, seis años planos. Ese
resultado no se obtiene de una proyección mecánica sino de un supuesto de
consolidación sostenida, y el capítulo 12 lo descompone.

\marginal{lo que el marco no trae}
Tres cosas que el marco no ofrece y que el resto del documento tiene que suplir o
declarar ausentes. No hay proyección de gasto por función, de modo que ninguna de
las trayectorias sectoriales de los capítulos 5 a 11 puede extenderse con material
oficial. No hay supuesto de población ni de matrícula, lo que deja sin denominador
a toda cifra per cápita. Y no hay una descomposición del cambio del acervo de
deuda en las variables que lo mueven, que es precisamente el hueco que el capítulo
12 llena con construcción propia.

\clearpage
